\section{The YASO Ledger Model}
\label{sec:ledger}

As discussed in Section~\ref{sec:overview-ledger}, the ledger combines a
synchronous public record with protected local computation.  We now use
it to define the execution of a YASO protocol.  The protocol supplies the
algorithms; the model fixes registration, communication, and corruption.
Our abstract model follows the role-based modeling approach of
Gentry et al.~\cite{GentryEtAl2021YOSOFull}; see Appendix~A of the full
version for a discussion of this choice.
All application parties, whether input providers or output recipients,
are subject to chosen corruption at any time.

\subsection{Protocol syntax}
\label{sec:yaso-syntax}

A YASO protocol is a tuple of uniform polynomial-time algorithms
\[
  \Pi=(\CRS,\Gen,\Schedule,\ComputationStep,\ApplicationStep).
\]
The security parameter $\lambda$ and session identifier $\sid$ are
implicit inputs where omitted.  The algorithms have the following interfaces.
\begin{description}
  \item[$\CRS(1^\lambda)\to\crs$.]
    Sample the public setup string.  The setup is trusted; its generation
    tape is not the state of a registered party.
  \item[$\Gen(1^\lambda,\crs,\Class)\to(\pk,\sk)$.]
    Generate public registration information $\pk$ and immutable
    private information $\sk$.  The class is either $\Application$ or
    $\Computation$.  We use $\pid=\pk$; $\pk$ may be a tuple of public
    keys whose interpretation belongs to $\Pi$.
  \item[$\Schedule(L)\to S$.]
    Deterministically return a list $S=(\pk_1,\ldots,\pk_k)$ of distinct
    registered public keys identifying the parties scheduled to speak.
    The list may be empty, $S=()$, when $L$ does not permit the protocol
    to progress.
  \item[$\ComputationStep(\pk,\sk,L)\to m$.]
    Compute a scheduled computation party's message
    $m$, or $\bot$ to remain silent.
  \item[$\ApplicationStep(\pk,\sk,L,x)\to(y,m)$.]
    Compute an application party's private output $y$ and message $m$.
    Either may be $\bot$.  The input $x$ is the party's ordered
    cumulative private input list, frozen at the start of the round.
    Each entry pairs an input with its arrival position in $L$.
\end{description}

The ledger records clock ticks as $\Tick$ entries in $L$; their number
is the current round $r$.  At the start of each round it fixes a snapshot
$L_r$, the list $S_r=\Schedule(L_r)$ of scheduled speakers, and an input
snapshot $x_{\pk,r}$ for each registered application party.
During the round, it executes $\ComputationStep$ for each scheduled honest
computation party and $\ApplicationStep$ for every honest application
party registered at round start, in an order chosen by the adversary.
Each executes at most once per round, using $L_r$ and, for application
parties, $x_{\pk,r}$.  Registrations, postings, and inputs received during
the round enter the snapshots for the following round.
The ledger stops executing a party's step
algorithm after it speaks.  Scheduling governs posting, not private
output: at the tick ending the round, an application party receives its
non-$\bot$ output whether or not it belongs to $S_r$.
A still-silent party may be scheduled again.

Each party's step algorithm uses a fixed private random tape $\rho_\pk$,
independent of its key-generation tape and read from the beginning on
each invocation.  Tape bits are sampled when first read and retained for
recomputation.  The ledger keeps the key-generation tape, all sampled
step-tape bits, inputs, and complete local computations in $\History_\pk$.
Private outputs are sent by $\FLedger[\Pi]$ to the application party
with identifier $\pid=\pk$.
The flag $\Spoken_\pk$ records whether the party has posted.

Public task data, including the next circuit and its input assignments,
are determined from $L$ by the protocol.  A party can propose such data
in its ledger message; the protocol determines which proposals to use.
The ledger does not interpret circuit descriptions or decide their validity.

\subsection{Ledger execution}
\label{sec:ledger-execution}

Fix $\Pi$, $\lambda$, $\sid$, and a computation-corruption probability
$p\in[0,1)$.  Figure~\ref{fig:yaso-ledger} defines
$\FLedger[\Pi]$; all commands and records are session-tagged.
We use the discrete-clock convention of GearBox~\cite{DavidEtAl2022GearBox}
(Section~2.1 of its full version): $\FClock$ supplies synchronized ticks,
and a time slot between successive ticks permits polynomial-time
computation.  Round $r$ occupies the slot starting at tick $r$, beginning
with $r=1$.

Within a round, the adversary requests individual executions through
$\Activate(\pk)$, interleaving them with postings and admissible
corruptions.  Computed private outputs are retained until the tick ending
the round.  Each tick after the first triggers an internal $\EndRound$
sweep: the ledger activates the remaining eligible honest parties,
records missing posts, and delivers the round's private outputs before
starting the next round.
Registration is available at any time.  Private application inputs may
arrive between activations and are included in the next round's input
snapshots.

We assume perfectly synchronized clocks and private output delivery at
the common tick.  In a model with bounded clock skew, parties would
instead output at their respective local ticks.

Corruptions refer to pids in the UC execution.  An application pid may be
corrupted at any time.  A computation pid is corrupted independently with
probability $p$ at registration; an initially honest computation pid may
be corrupted only after it has spoken.  The ledger's corruption flags
record this same status, and an allowed corruption discloses the retained
state associated with that pid.  The ledger neither checks signatures nor
interprets posted messages.

\begin{figure}[p]
\centering
\fbox{%
\begin{minipage}{0.95\textwidth}
\small
\textbf{The ledger $\FLedger[\Pi]$.}
Public state consists of $L$; $r$ is the number of $\Tick$ entries in $L$.
For each registered $\pk$, retain its class, $\sk_\pk$, random tape
$\rho_\pk$, full history $\History_\pk$, cumulative input list $x_\pk$, and flags
$\Spoken_\pk,\Corrupted_\pk$.
After the first tick, also retain the fixed snapshot $L_r$ and input
snapshots $x_{\pk,r}$ for application pids registered at round start,
a list $S_r$ of speakers, a list $E_r$ of eligible pids, and a set $A_r$
of pids already activated in round $r$.  Retain the round's pending
private outputs in a list $\PendingOutputs_r$.

\begin{description}
\item[Initialization.]
  Sample $\crs\sample\CRS(1^\lambda)$, set
  $L=((\Setup,\sid,\crs))$ and initialize an empty registry.

\item[$\Register(\Class)$ from the adversary.]
  For a valid class, sample $(\pk,\sk_\pk)\sample
  \Gen(1^\lambda,\crs,\Class)$.
  Reject an already registered $\pk$; otherwise store its class and
  $\sk_\pk$.  Initialize $\History_\pk$ with $\Gen$'s input tuple, random tape,
  and sequence of local configurations through termination, including
  its returned keys.  Initialize an independent private
  step tape $\rho_\pk$, set $x_\pk=()$ and both flags false, and append
  $(\Register,\pk,\Class)$ to $L$.
  For class $\Computation$, independently set
  $\Corrupted_\pk\sample\mathsf{Bernoulli}(p)$ and, if true, disclose
  its state and history to the adversary; pid $\pk$ is then corrupted
  from registration.  Return $\pk$ to the adversary.  For class
  $\Application$, also output $\pk$ to pid $\pk$.

\item[$\Read$ from any pid or from the adversary.]
  Output the current $L$ to the requesting pid or adversary, respectively.

\item[$\Input(u)$ from party $\pid=\pk$.]
  If $\pk$ is an honest application party, append $(|L|,u)$ to $x_\pk$
  and record its arrival in $\History_\pk$.
  Ignore inputs to computation parties; route a corrupted
  application party's inputs to the adversary.

\item[Clock tick from $\FClock$.]
  Before the first tick, ignore $\Activate$ and $\Post$ commands.
  If $r>0$, perform the internal $\EndRound$ sweep below for round $r$.
  Append $\Tick$ to $L$, let $r$ be its tick count, and fix $L_r=L$.
  For each registered application pid $\pk$, fix $x_{\pk,r}=x_\pk$.
  Compute $S_r=\Schedule(L_r)$ and discard unknown keys and parties
  that have already spoken.
  Let $E_r$ list all unspoken application pids and all unspoken computation
  pids in $S_r$, in registration order.  Set $A_r=\varnothing$ and
  $\PendingOutputs_r=()$.

\item[$\Activate(\pk)$ from the adversary.]
  Require $\pk\in E_r\setminus A_r$ and $\Corrupted_\pk=\bot$.
  If $\Class_\pk=\Computation$, run
  \[
    m_\pk\gets\ComputationStep(\pk,\sk_\pk,L_r;\rho_\pk).
  \]
  Otherwise run
  \[
    (y_\pk,m_\pk)\gets\ApplicationStep(\pk,\sk_\pk,L_r,x_{\pk,r};\rho_\pk).
  \]
  Append to $\History_\pk$ the round $r$, the input tuple supplied to
  the step algorithm, the random bits read, and its sequence of local
  configurations through termination, including its returned value.
  Add $\pk$ to $A_r$.
  If $\pk\in S_r$ and $m_\pk\neq\bot$, append
  $(\Message,r,\pk,m_\pk)$ to $L$ and set $\Spoken_\pk=\top$.
  For an application pid with $y_\pk\neq\bot$, append $(\pk,y_\pk)$
  to $\PendingOutputs_r$.

\item[$\Post(\pk,m)$ from the adversary.]
  Require $m\in\{0,1\}^*$, $\pk\in S_r$, $\Corrupted_\pk=\top$,
  and $\Spoken_\pk=\bot$.
  Append $(\Message,r,\pk,m)$ to $L$ and set $\Spoken_\pk=\top$.

\item[$\EndRound$ (internal sweep at a clock tick).]
  Visit $E_r$ in registration order.  For each $\pk$ still honest and
  not in $A_r$, perform the steps of $\Activate(\pk)$.
  Append $(\NoMsg,r,\pk)$ to $L$ for every $\pk\in S_r$ that did not
  post, in registration order.
  For each $(\pk,y)\in\PendingOutputs_r$, in registration order,
  output $y$ to pid $\pk$.

\item[On corruption of pid $\pk$.]
  The adversary may corrupt this pid if $\Class_\pk=\Application$ or
  $\Spoken_\pk=\top$.  Set $\Corrupted_\pk=\top$ and disclose its
  $\sk_\pk$, all sampled bits of $\rho_\pk$, $x_\pk$, its speaking
  flag, and complete retained history to the adversary.
  Subsequently the ledger no longer executes the pid's step algorithm.
  The adversary can use $\Post$ to post on behalf of the corrupted pid.
\end{description}
Reject malformed commands and commands naming an
unregistered key, except that $\Read$ is available
to any pid.
\end{minipage}}
\caption{Execution of a YASO protocol.  Silence preserves protection for
an initially honest computation party; only an actual posting sets
$\Spoken$.}
\label{fig:yaso-ledger}
\end{figure}

The parameter $p$ belongs to the execution model.  Committee sizes and
sharing thresholds belong to the concrete protocol, not to the ledger.
Security is stated separately by specifying a target functionality
$\mathcal F$ and requiring a simulation of this execution in the UC
framework.  The ongoing $\FMPC$ sketched in
Section~\ref{sec:overview-functionality} is our intended target; neither
it nor its circuit descriptions are additional algorithms in $\Pi$.
